%% =============================================================================
%% sm_f7_decomposition.tex — SM-E, Subsection E.7: Coverage Gap Decomposition
%% Body content only (\subsection level); parent structure in sm_e.tex
%% =============================================================================

\subsection{Coverage gap decomposition}
\label{smf:decomposition}

\Cref{smf:tab-decomposition} decomposes the coverage gap for each
parameter--estimator combination under the S3 (NSECE-calibrated) scenario.
Two diagnostic quantities distinguish width-driven from bias-driven
shortfalls:
\begin{itemize}[nosep]
  \item \emph{SE calibration ratio} $= \widehat{\mathrm{SD}} /
    \overline{\mathrm{SE}}$, where $\widehat{\mathrm{SD}}$ is the
    empirical standard deviation of point estimates across $R = 200$
    replications and $\overline{\mathrm{SE}}$ is the mean reported
    standard error.  A ratio exceeding 1 indicates that the reported SE
    underestimates the actual sampling variability (width-driven gap).
  \item \emph{Standardized mean bias} $= |\bar{b}| /
    \widehat{\mathrm{SD}}$, where $\bar{b}$ is the mean bias across
    replications.  Values exceeding $\sim 0.5$ indicate non-negligible
    centering error (bias-driven gap).
\end{itemize}

\begin{table}[htbp]
\centering
\caption{Coverage gap decomposition under the S3 (NSECE-calibrated)
  scenario.  ``Below'' and ``Above'' report the percentage of replications
  where the true value falls below or above the confidence interval,
  respectively.  SE ratio $> 1$ indicates intervals that are too narrow;
  $|\bar{b}|/\mathrm{SD} > 0.5$ indicates non-negligible point estimate
  bias.  $\star$ = sandwich structurally inapplicable (E-WS $\equiv$ E-WT).}
\label{smf:tab-decomposition}
% Generated by 87_B8_coverage_decomposition.R — S3 (NSECE-calibrated) scenario
\small
\begin{adjustbox}{max width=\textwidth}
\begin{tabular}{@{}l l r rr r r l@{}}
\toprule
 & & & \multicolumn{2}{c}{Non-coverage (\%)} & & \\
\cmidrule(lr){4-5}
Parameter & Est. & Cov.\ (\%) & Below & Above &
  SE ratio & $|\bar{b}|/\text{SD}$ & Dominant source \\
\midrule
$\alpha_{\text{pov}}$ & E-UW & 89.5 & 1.0 & 9.5 & 0.86 & 0.63 & ---  \\
 & E-WT & 58.5 & 15.5 & 26.0 & 1.93 & 0.22 & Width \\
 & E-WS & 82.0 & 5.0 & 13.0 & 1.13 & 0.22 & Width \\[3pt]
$\beta_{\text{pov}}$ & E-UW & 89.0 & 0.5 & 10.5 & 0.85 & 0.69 & --- \\
 & E-WT & 62.5 & 15.5 & 22.0 & 1.88 & 0.06 & Width \\
 & E-WS & 82.5 & 8.5 & 9.0 & 1.14 & 0.06 & Width \\[3pt]
$\log\kappa$ & E-UW & 92.0 & 1.0 & 7.0 & 0.85 & 0.48 & --- \\
 & E-WT & 34.0 & 65.0 & 1.0 & 1.75 & 1.37 & Both \\
 & E-WS & 60.5 & 39.5 & 0.0 & 1.00 & 1.37 & Bias \\[3pt]
$\tau_{\text{ext}}$ & E-UW & 98.0 & 1.5 & 0.5 & 0.64 & 0.59 & --- \\
 & E-WT/WS & 0.5 & 99.5 & 0.0 & 1.22 & 3.18 & Bias$^{\star}$ \\[3pt]
$\tau_{\text{int}}$ & E-UW & 97.5 & 1.0 & 1.5 & 0.73 & 0.51 & --- \\
 & E-WT/WS & 0.0 & 100.0 & 0.0 & 1.11 & 3.55 & Bias$^{\star}$ \\
\bottomrule
\end{tabular}
\end{adjustbox}

\end{table}

For the poverty coefficients under E-WS, the SE ratio of 1.13--1.14
and negligible standardized bias (0.06--0.22) confirm that the 8~pp
coverage gap is entirely width-driven: the sandwich variance estimator
underestimates the true sampling variance by approximately 13\%.
This is consistent with the well-documented $O(S^{-1})$ finite-sample
downward bias of cluster-robust variance estimators when $S = 51$
\citep{Binder1983}.  For $\log\kappa$, the picture reverses: the SE
ratio is 1.00 (perfectly calibrated width), but the standardized bias
of 1.37 indicates that the pseudo-likelihood point estimate is
systematically shifted, producing entirely one-sided non-coverage
(39.5\% below, 0\% above).  For $\tau_{\mathrm{ext}}$ and
$\tau_{\mathrm{int}}$, massive standardized bias ($> 3$) overwhelms
any width consideration, reaffirming the structural inapplicability of
the sandwich correction to variance components.
